\section{Disegno sperimentale}\label{sec:exp}

Il protocollo è articolato in tre studi. Lo Studio~I seleziona il passo per la configurazione di riferimento e lo verifica su repliche indipendenti. Lo Studio~II misura la scalabilità e confronta le modalità sincrona e asincrona. Lo Studio~III valuta la competitività del CPSO rispetto alle varianti PSO e al controllo discreto, con budget di valutazioni e durate prefissati. La qualità a budget di valutazioni fissato caratterizza la dinamica della ricerca. La qualità entro una scadenza include anche il costo di esecuzione.

Le misure sono raccolte su un MacBook Air con Apple M2, dotato di quattro core ad alte prestazioni e quattro core ad alta efficienza \citep{AppleM2Specs,AppleM2Design}. Negli Studi I e II si usano $N=64$ particelle e coefficienti casuali scalari. Per il CPSO si fissano $c_c=c_s=2.05$ e inerzia lineare da $0.9$ a $0.4$. Le posizioni sono inizializzate uniformemente nel dominio e le velocità in $[-0.1w_j,0.1w_j]$ per coordinata, dove $w_j$ è l'ampiezza del dominio. Una particella che oltrepassa un limite è riportata sul bordo e la componente di velocità uscente è invertita e dimezzata. Le funzioni obiettivo sono definite in Appendice~\ref{app:dynamics}.

\subsection{Sensibilità al passo}\label{subsec:exp-dt}
Si confrontano 24 passi in $[0.01,50]$ sulle sei funzioni standard, a $D=30,100$. Ciascun budget, $2\,048$, $8\,192$ o $32\,768$ FE, corrisponde a un'esecuzione indipendente con il proprio profilo d'inerzia. Misuriamo la riduzione del residuo rispetto al migliore valore iniziale:
\begin{equation}
 z(h;B)=\log_{10}\!\left[\operatorname{clip}_{[10^{-16},1]}
 \frac{\max\{J^{(B)}(h)-J^*,0\}}{J^{(0)}-J^*}\right].
 \label{eq:dt-criterion}
\end{equation}
Il passo che minimizza la media di $z$ su funzioni, dimensioni e budget è scelto su 30 repliche appaiate. La verifica usa 30 repliche indipendenti e confronta il passo scelto con $dt=0.35$, $0.5$ e $30$. Gli intervalli al 95\% sono ottenuti mediante bootstrap delle repliche appaiate. Ai contrasti multipli si applica la correzione di Holm.

\subsection{Scalabilità e costo della valutazione}\label{subsec:exp-parallel}
Lo Studio II presenta il sottoinsieme delle misure della campagna relativo alla diffusione inversa. Si usa il passo selezionato nello Studio I. Il confronto esplorativo tra gli otto vicinati impiega il particle-parallel sincrono a $D=30,100$ e usa come riferimento il tempo a un worker per ciascun grafo. La verifica indipendente comprende 30 semi osservati in due sessioni di misura. A $8\,192$ FE, con vicinato Global, si confrontano la seriale e i backend master--slave e particle-parallel a $p=1,2,4,8$ worker, nelle stesse due dimensioni e con 16, 64, 256 o 512 passi RK4 per valutazione. A $32\,768$ FE si usa la risoluzione di 512 passi e si confrontano la seriale e i due backend a quattro e otto worker. A $D=30$ si ripete quest'ultimo confronto con Random-8.

Per un esecutore sincrono $a$, accelerazione rispetto alla seriale, strong scaling ed efficienza sono
\begin{equation}
 A_a(p)=\frac{T_s}{T_{a,p}},\qquad
 S_a(p)=\frac{T_{a,1}}{T_{a,p}},\qquad E_a(p)=\frac{S_a(p)}p.
 \label{eq:speedup-efficiency}
\end{equation}
Il rapporto $T_{a,4}/T_{a,8}$ misura il guadagno incrementale a parità di grafo, budget e discretizzazione. Le accelerazioni e i rapporti confrontano i tempi nativi di esecuzioni complete. Nella verifica indipendente si calcola la media geometrica dei rapporti appaiati e si ottengono intervalli bootstrap puntuali al 95\%, aggregando le due sessioni per seme prima del ricampionamento. I test sul rapporto fra quattro e otto worker conservano la correzione di Holm applicata al blocco completo della campagna. A parità di FE si controllano il valore iniziale, il miglior valore finale e la traccia di ricerca rispetto alla seriale abbinata.

L'esecuzione asincrona è confrontata con il particle-parallel a $8\,192$ FE e 512 passi RK4, con quattro e otto worker. Il confronto usa Global a $D=30,100$ e Random-8 a $D=30$. Si misurano il rapporto dei tempi $T_{\mathrm{sync}}/T_{\mathrm{async}}$ e la differenza di qualità $z_{\mathrm{async}}-z_{\mathrm{sync}}$. Le tre ripetizioni asincrone per replica e sessione sono aggregate prima del ricampionamento.

\subsection{Confronto con varianti PSO}\label{subsec:exp-pso}
La configurazione CPSO è selezionata con un fattoriale $2\times2\times2$. Si incrociano $dt\in\{0.38125,1\}$, i vicinati Global e Random-8, e la casualità scalare o per coordinata. Ogni configurazione usa 100 particelle e quattro worker. La selezione minimizza il rango medio del gap finale sui blocchi appaiati di Rastrigin, Rosenbrock e Ackley, a $D=20,100$, con $50\,000$ e $500\,000$ FE e cinque semi distinti da quelli della valutazione successiva. Un secondo blocco con cinque ulteriori semi mantiene fissi vicinato, casualità e passo di base. L'adattamento del passo, il recupero di particelle e la loro combinazione sono confrontati con la configurazione fissa. La regola di selezione richiede una riduzione media del log-gap di almeno $0.05$ per adottare un adattamento. Il riferimento CPSO Global con casualità scalare e $dt=0.38125$ misura l'effetto congiunto delle scelte di ricerca.

Il confronto a valutazioni fissate include la PSO discreta abbinata e dieci varianti PSO di letteratura (Appendice~\ref{app:baselines}). Il controllo abbinato condivide con il CPSO selezionato popolazione, inizializzazione, vicinato, semi, estrazioni per coordinata, inerzia, coefficienti di accelerazione, gestione dei bordi e numero di worker. A ogni generazione sostituisce la soluzione analitica con l'aggiornamento discreto a passo unitario delle Eq.~\eqref{eq:or-discrete-position} e~\eqref{eq:or-discrete-velocity}, usando il vicinato Random-8 al posto dell'attrattore globale. Le due dinamiche usano la stessa proiezione al bordo con inversione e dimezzamento della componente di velocità uscente. Il controllo discreto usa velocità senza saturazione, $c_c=c_s=2.05$, inizializzazione uniforme in $[-0.1w_j,0.1w_j]$ per componente e inerzia lineare da $0.9$ a $0.4$ sul budget di aggiornamenti.

I confronti a FE fissate impiegano $N=100$ particelle per metodo. Le sei funzioni standard si valutano a $D=20,100$, con 20 repliche a semi distinti da quelli di sviluppo e budget indipendenti di $50\,000$ e $500\,000$ FE. Rastrigin, Rosenbrock e Ackley sono comuni a selezione e confronto nativo, che usa le repliche aggiuntive per valutare la ripetibilità. La verifica esterna comprende le funzioni COCO BBOB $F_{2}$, $F_{4}$, $F_{5}$, $F_{7}$, $F_{9}$, $F_{11}$, $F_{13}$, $F_{14}$, $F_{17}$, $F_{18}$, $F_{19}$, $F_{20}$, $F_{22}$, $F_{24}$, con cinque istanze (11--15), a $D=20,40$ e budget di $10\,000$ e $50\,000$ FE \citep{Hansen2021,Hansen2016Procedure}.

Il confronto a tempo fissato usa Rastrigin e diffusione inversa (512 passi RK4), a $D=20,100$ con 20 repliche. Si assegnano 15 o 60 s al CPSO a uno e quattro worker e a CSO, CLPSO, SL-PSO e SPSO-2011 a quattro worker. Tutti usano lo stesso valutatore e il risultato è il migliore valore disponibile al termine dell'ultimo lotto completato entro la scadenza. I profili d'inerzia dipendenti dal numero di iterazioni usano un orizzonte comune di $1\,000$ aggiornamenti. Nel CPSO il peso scende linearmente da $0.9$ a $0.4$ e rimane a $0.4$ dopo tale orizzonte, per entrambe le configurazioni di worker. Per il target $g\leq10^{-8}$ si riportano la frazione di successi entro la scadenza e la mediana dei tempi fra i successi. Il confronto interno al CPSO verifica che la traccia del migliore valore coincida sul prefisso comune di FE.

La qualità finale è il gap $g=\max(J-J^*,10^{-12})$, dove $J^*$ è il valore ottimo noto per la funzione o l’istanza. I ranghi sono calcolati nei blocchi appaiati e aggregati dando uguale peso alle funzioni. Ogni contrasto per funzione è la differenza media appaiata di $\log_{10}(g)$ sulle repliche o sulle istanze. Valori entro $\pm0.01$ definiscono la parità pratica. Il confronto con la migliore variante di letteratura considera, per ciascuna funzione, il metodo con il log-gap medio più basso. I test bilaterali dei segni usano le funzioni come unità e la correzione di Holm entro ciascuna campagna. Per BBOB la famiglia comprende $48$ contrasti, uno per ciascuna combinazione di $12$ avversari, due dimensioni e due budget. Le ECDF misurano la frazione di coppie esecuzione--target raggiunte per i livelli $10^2,1,10^{-2},10^{-3},10^{-4},10^{-5},10^{-6},10^{-8}$, mantenendo nel denominatore anche le coppie che non raggiungono il target.
